\section{Ablation Study}
\label{sec:ablations}

We conduct ablation experiments to quantify the contribution of each
feature family. All ablations use the same LightGBM hyperparameters
(\texttt{n\_estimators}${}=500$, lr${}=0.05$,
\texttt{min\_samples\_leaf}${}=5$) and are evaluated on the frozen
expert test set ($n=400$).

\begin{table}[t]
\centering
\caption{Feature-family ablation on \texttt{human\_test\_frozen}
($n=400$). All variants use the same LightGBM hyperparameters.}
\label{tab:ablation}
\begin{tabular}{lrrrr}
\toprule
Feature set & Dims & Tier Acc & Macro F1 & Adj Acc \\
\midrule
LLM only & 7 & 0.468 & 0.418 & 0.858 \\
Structural only & 13 & 0.393 & 0.317 & 0.750 \\
Opus only & 2 & 0.523 & 0.382 & 0.843 \\
LLM + Struct & 20 & 0.495 & 0.444 & 0.855 \\
LLM + Opus & 9 & 0.463 & 0.416 & 0.818 \\
Struct + Opus & 15 & 0.495 & 0.432 & 0.835 \\
\midrule
\textbf{All (LLM+Struct+Opus)} & \textbf{22} & \textbf{0.533} & \textbf{0.489} & \textbf{0.843} \\
\bottomrule
\end{tabular}
\end{table}


\begin{figure}[t]
\centering
\includegraphics[width=\linewidth]{fig_feature_importance}
\caption{LightGBM feature importance. (A)~Per-feature importance for
all 22 features, colored by family. (B)~Cumulative importance curve.
The top five features capture over 60\% of the total importance budget.}
\label{fig:feature_importance}
\end{figure}

\begin{figure}[t]
\centering
\includegraphics[width=\linewidth]{fig_ablation_bar}
\caption{Feature family ablation results on the frozen test set
($n=400$). Tier accuracy for GBM models trained with different
feature-family subsets. The full model (All) achieves the highest
accuracy.}
\label{fig:ablation_bar}
\end{figure}

\subsection{Key Findings}

\begin{itemize}
    \item \textbf{All three families contribute.} The full 22-feature
          model (0.533 accuracy, 0.489 macro F1) outperforms every
          subset, confirming that the three signal families capture
          complementary information.

    \item \textbf{Opus scores are the strongest single channel.} The
          two-dimensional Opus-only model achieves 0.523 accuracy,
          nearly matching the full model. This demonstrates that a
          well-calibrated LLM-as-judge provides a strong prior that the
          stacker refines rather than replaces.

    \item \textbf{LLM features outperform structural features in
          isolation.} LLM-only (7d) reaches 0.468 accuracy versus 0.393
          for Structural-only (13d), a 7.5 percentage-point gap. However,
          Structural features provide the largest marginal gain when
          combined: LLM+Struct (0.495) exceeds LLM-only by 2.7~pp,
          while LLM+Opus (0.463) is slightly below LLM-only, suggesting
          the Opus signal partially overlaps with the LLM embedding
          signal.

    \item \textbf{Structural features reduce adjacent-accuracy
          degradation.} Structural-only has the lowest adjacent accuracy
          (0.750), but adding structural features to Opus (Struct+Opus:
          0.835) provides better ordinal calibration than LLM+Opus
          (0.818), indicating that graph topology helps the model
          respect the ordinal tier structure.

    \item \textbf{The full model is Pareto-optimal.} No subset achieves
          higher accuracy, F1, or adjacent accuracy than the full model,
          justifying the inclusion of all three feature families despite
          the modest dimensionality overhead.
\end{itemize}

These results confirm that the stacking approach, combining text
similarity, graph structure, and LLM judgment in a single gradient
boosting model, captures relational structure that no single signal
family provides alone.

\subsection{Loss Function Comparison}

The v\_final pipeline trains three encoder variants, one per loss
function, to assess the tradeoff between ordinal supervision strength
and discriminative sharpness:

\begin{itemize}
    \item \textbf{KL ordinal vs.\ Focal:} KL ordinal loss injects
          explicit Gaussian smoothing over adjacent tiers, which is
          expected to improve ordinal calibration (adjacent accuracy)
          at the cost of some peak-tier precision. Focal loss retains
          hard targets but concentrates gradient on difficult boundary
          pairs, trading ordinal smoothness for sharper discrimination
          at class boundaries. Preliminary expectation is that KL
          ordinal achieves higher adjacent accuracy while Focal achieves
          higher binary accuracy.

    \item \textbf{CORN vs.\ KL ordinal:} CORN decomposes the task
          into $K{-}1$ binary problems, guaranteeing monotone
          cumulative probability estimates by construction. This hard
          constraint may reduce flexibility for complex distributions
          but prevents the probability inversions (e.g., $P(\text{Related})
          > P(\text{Partial})$ with $P(\text{Unrelated})$ near zero)
          occasionally produced by softmax outputs under KL training.

    \item \textbf{Ensemble benefit:} Because the three losses produce
          diverse probability estimates, the meta-stacker benefits from
          complementary error patterns. Early experiments suggest
          ensemble gains of 2--4 percentage points in macro F1 relative
          to the best single-loss model. Actual figures will be reported
          once training completes.
\end{itemize}

\begin{figure}[t]
\centering
% \includegraphics[width=\linewidth]{fig_loss_comparison}
\fbox{\parbox{\linewidth}{\centering\vspace{2em}
Loss function comparison figure --- to be generated after training\vspace{2em}}}
\caption{Per-class F1 for each loss function (KL ordinal, CORN, Focal)
on the frozen test set. Bars are grouped by tier; error bars show
bootstrap 95\% CIs. The ensemble (all three) is shown as a separate
group.}
\label{fig:loss_comparison}
\end{figure}

\subsection{Architecture Comparison}

The three encoder architectures span the cross-encoder / bi-encoder
design space:

\begin{itemize}
    \item \textbf{Cross-encoder accuracy vs.\ bi-encoder speed:}
          RoBERTa-large and DeBERTa-base process the full concatenated
          pair through joint self-attention, giving the model access to
          fine-grained token-level interactions. BGE-large-v1.5 encodes
          each text independently, enabling offline pre-computation of
          embeddings and $O(1)$ inference per query. The expected
          accuracy ordering is RoBERTa-large $\geq$ DeBERTa-base
          $>$ BGE-large-v1.5, but the bi-encoder CLS similarity features
          contribute distinct signal to the meta-stacker that neither
          cross-encoder alone provides.

    \item \textbf{DeBERTa-base vs.\ RoBERTa-large:} DeBERTa's
          disentangled attention over content and position embeddings
          confers a representational advantage on structurally regular
          inputs (e.g., bullet-pointed control descriptions), potentially
          closing the gap with the larger RoBERTa model on this corpus.

    \item \textbf{Meta-stacker feature attribution:} Coefficient
          magnitudes in the L2-regularized logistic regression
          meta-stacker will be reported after training and are expected
          to replicate the v7c finding that the largest-capacity
          cross-encoder (here RoBERTa-large) contributes the most
          influential logit features, with the bi-encoder cosine features
          providing meaningful complementary signal.
\end{itemize}

\begin{figure}[t]
\centering
% \includegraphics[width=\linewidth]{fig_arch_comparison}
\fbox{\parbox{\linewidth}{\centering\vspace{2em}
Architecture comparison figure --- to be generated after training\vspace{2em}}}
\caption{Speed--accuracy tradeoff for the three v\_final encoder
architectures. X-axis: median inference latency per pair (ms, log
scale). Y-axis: macro F1 on the frozen test set. The bi-encoder
(BGE-large-v1.5) is orders of magnitude faster at a modest accuracy
cost; the ensemble stacker achieves the best accuracy.}
\label{fig:arch_comparison}
\end{figure}

\subsection{Validation Split Impact}

Mapping-level deduplication fundamentally changes the difficulty of
the validation task relative to the v7c pair-level split:

\begin{itemize}
    \item \textbf{Contamination rate before dedup:} 56\% of validation
          text pairs had near-duplicate counterparts in the training set
          under the v7c pair-level split strategy. This inflated
          validation metrics and obscured genuine generalization gaps.

    \item \textbf{After mapping-level dedup:} All pairs sharing the same
          source--target CRE mapping are assigned exclusively to either
          the training or validation split. This eliminates the
          possibility of a model memorizing a specific pair during
          training and recovering it in validation with minor surface
          variation.

    \item \textbf{Expected metric shift:} Mapping-level dedup is expected
          to lower raw validation accuracy by approximately 3--7
          percentage points relative to the pair-level baseline,
          reflecting a more honest estimate of generalization. Models
          that close this gap through architecture and loss improvements
          rather than memorization demonstrate genuine transfer.

    \item \textbf{Impact on CollapseGuard and OverfittingGuard:} A
          harder, non-contaminated validation set provides more reliable
          early-stopping signals to the training guards. False ``good''
          checkpoints induced by memorized validation pairs in v7c
          training are eliminated, making the guards' rollback decisions
          more accurate.
\end{itemize}
